\section{三种方法的统一比较}

\begin{table}[H]
\centering
\small
\begin{tabular}{>{\raggedright\arraybackslash}p{2.2cm}>{\raggedright\arraybackslash}p{3.4cm}>{\raggedright\arraybackslash}p{3.4cm}>{\raggedright\arraybackslash}p{4.2cm}}
\toprule
\textbf{方法} & \textbf{主要阶段} & \textbf{核心表示} & \textbf{主要收益 / 代价} \\
\midrule
LATS & 推理时搜索 & 状态--动作树 & 可回溯、探索多计划；调用成本高 \\
SPRINT & 训练后并行推理 & 任务依赖 DAG & 降低关键路径；依赖识别和汇总复杂 \\
SWiRL & 离线数据 + RL & 多步工具轨迹 & 学到通用过程；受 teacher、judge 与离线覆盖限制 \\
\bottomrule
\end{tabular}
\end{table}

\subsection{组合方式}

三者并不互斥。一个完整系统可以用 SPRINT 生成并行任务图，在每个困难子任务内运行 LATS 搜索，再用 SWiRL 训练得到的 policy 选择工具与处理 observation。

\subsection{工程选择}

\begin{itemize}
    \item 任务价值高、可模拟、需回溯：优先树搜索；
    \item 子任务独立、工具延迟高：优先并行执行；
    \item 有大量可验证环境、希望长期改进：构造 synthetic trajectories 做 RL；
    \item 动作不可逆或环境昂贵：限制在线探索，增加离线预演和人工审批。
\end{itemize}

\subsection{本章小结}

规划问题可以从搜索、调度和学习三个层次解决。树搜索提高决策质量，并行调度降低时间，过程级训练把经验沉淀进 policy。三者共同目标是减少长轨迹中的错误累积与无效计算。

